\section{その他}

\subsection{企業年金連合会に関する通則}

出題例なし

\newpage

\subsection{企業年金連合会の設立及び管理}

出題例なし

\newpage

\subsection{企業年金連合会の行う業務}

出題例なし

\newpage

\subsection{企業年金連合会からDB、DCへの積立金の移換}

出題例なし

\newpage

\subsection{企業年金連合会の解散及び清算}

出題例なし

\newpage

\subsection{税制上の措置}

税制の問題集に記載。

\newpage

\subsection{業務の委託}

出題例なし

\newpage

\subsection{事業年度その他財務}

\begin{truefalse}[年金2 2020 問題1(1)\textcircled{2}]
  下線部の正誤を判定し、誤っている場合は正しい内容に改めなさい。
  \\

  確定給付企業年金を実施する企業年金基金の前事業年度の末日の年金経理において、純資
  産額が 100 億円、数理上資産額が 105 億円、数理債務が 100 億円、責任準備金が
  95 億円、 最低積立基準額が 85 億円である場合、\underline{5 億円}を限度として、年金経理から業務経理へ繰り入
  れることができる。なお、年金経理から業務経理へ繰り入れるための要件は満たしているものとする。

\end{truefalse}

\begin{answerbox}
  $数理上資産額 - 責任準備金 = 105 - 95 = 10億円 \cdots \textcircled{1}$

  $純資産額 - 最低積立基準額 = 100 - 85 = 15億円 \cdots \textcircled{2}$

  \textcircled{1}と\textcircled{2}のうち低い額を限度に年金経理から業務経理へ繰り入れることができるので、10億円が限度。
\end{answerbox}


\begin{shortanswer}[数理人試験 2025 問題4 設問2]
  確定給付企業年金法施行規則第百十三条に定める事業年度を1年としないことができる場合について簡記せよ。
\end{shortanswer}


\begin{answerbox}
  \begin{enumerate}[1.]
    \item 事業年度を変更した場合
    \item \begin{itemize}
      \item 法第三条第一項の規定により確定給付企業年金を実施しようとする場合
      \item 法第七十四条第一項の規定により規約型企業年金を他の規約型企業年金
            と統合する場合、法第七十五条第一項の規定により規約型企業年金を分割する場合、法第七十六
            条第三項若しくは法第七十七条第四項の規定により合併若しくは分割によって基金を
            設立する場合又は法第八十条第二項若しくは法第八十一条第二項の規定により給付の
            支給に関する権利義務を承継する場合（規約型企業年金を実施することとなる場合又
            は基金を設立することとなる場合であって、給付の支給に関する権利義務の承継に係
            る確定給付企業年金の掛金の額を給付の支給に関する権利義務の移転に係る確定給付
            企業年金の掛金の額と異なるものとする場合に限る。）
      \item 法第五十八条第一項の規定により財政再計算を行う場合
    \end{itemize}
  \end{enumerate}
\end{answerbox}

\newpage

\subsection{年金数理、年金数理人による確認}

出題例なし

\newpage

\subsection{書類等の提出・届出、報告}

\begin{fillblank}[数理人試験 2021 問題4 設問1]
  \noindent \textbf{◯承認認可基準別紙2}
  \begin{lawstructure}
    \begin{lawarticle}{}{}
      \lawparagraph{\textcircled{3}}
      監事の監査は、別紙5の「企業年金基金監事監査規程要綱」を基準として 監査規定を設け、これに基づき適正かつ厳正に行うこと。また、総合型企業 年金基金にあっては、規則第 117 条第4項第1号及び第2号並びに確定給付 企業年金法施行規則第 117 条第4項第3号に規定する監査に準ずるものとし て厚生労働大臣が定めるもの(厚生労働省告示第335号)の規定に基づき、 常時（　　A　　）以上の積立金を積み立て、又は積み立てると見込まれる場合は、 当該要件に該当した決算の翌々年度決算から、次のいずれかを受け、その結 果を監事の監査に活用して監事の監査の充実を図ること。

      \lawitem{ア}
      （　　B　　）又は公認会計士による会計監査

      \lawitem{イ}
      別紙5の2の「（　　C　　）のチェック項目及びチェックポイント」
      及び日本公認会計士協会業種別委員会実務指針第 62 号「総合型確定給付企 業年金基金に対する（　　C　　）業務に
      関する実務指針」(以下「実務指針」という。)等に基づき公認会計士又は（　　B　　）(これらの者と同等水準
      で業務を遂行できる者を含み、当該基金の理事及び職員を除く。)とあらかじめ手続を合意し、
      監事の監査に帯同する等して実施する（　　C　　）

      （以下略）
    \end{lawarticle}
  \end{lawstructure}
\end{fillblank}

\begin{answerbox}
  \begin{enumerate}[A: ]
    \item 20億円
    \item 監査法人
    \item 合意された手続き
  \end{enumerate}
\end{answerbox}

\begin{truefalse}[年金1 2018 問題1(1)\textcircled{3}]
  下線部の正誤を判定し、誤っている場合は正しい内容に改めなさい。
  \\ 

  総合型の企業年金基金は、貸借対照表（年金経理）の資産総額が\underline{30 億円を超えた決算の翌
年度}決算から、会計監査か合意された手続のいずれかを受け、その結果を監事の監査に活用
して監事の監査の充実を図ることとされている。
\end{truefalse}

\begin{answerbox}
  \textbf{誤}

  20億円を超えた決算の翌々年度
\end{answerbox}

\begin{explanationbox}
  承認認可基準別紙2(4)\textcircled{3}
\end{explanationbox}

\begin{shortanswer}[年金1 2022 問題2(1)エ]
  確定給付企業年金法施行規則第117 条第4 項において、企業年金基金が事業及び決算に関す
る報告書を地方厚生局長等に提出する場合、2 以上の事業主が共同して設立する企業年金基金
（※）は、会計監査等の結果を考慮した監事の意見を付けて代議員会に提出し、その議決を
得なければならないことと定められている。
\\

（※）同規則第19 条の2 第1 号に掲げる要件に該当する企業年金基金及び積立金の額が常時
20 億円を下回る、又は下回ると見込まれる企業年金基金を除く。
\\

この「積立金の額が常時20 億円を下回る、又は下回ると見込まれる」こととして、 『確定給
付企業年金法並びにこれに基づく政令及び省令について（法令解釈）』に記載されている内容
を簡潔に入力しなさい。
\end{shortanswer}

\begin{answerbox}
  次の(1)又は(2)のいずれかに該当することであること

  \begin{enumerate}[(1)]
    \item 直近３ヵ年の決算のうち少なくとも２ヵ年において、積立金（純資産）の額が20億円未満であること
    \item 直近３ヵ年の決算のうち少なくとも２ヵ年において、責任準備金の額又は最低積立基準額のいずれか低い額が20億円未満であること
  \end{enumerate}
  
\end{answerbox}

\begin{truefalse}
  下線部の正誤を判定し、誤っている場合は正しい内容に改めなさい。
  \\   

  二以上の事業主が共同して設立する基金（確定給付企業年金法施行規則第十九条の二第一号に
掲げる要件に該当する基金及び積立金の額が\underline{常時百億円を下回る基金}を除く。）の監事である
場合にあっては、監査法人の監査の結果、公認会計士の資格を有する者の監査の結果またはそ
れらに準ずるものとして厚生労働大臣が定めるものの結果のいずれかを考慮した意見を確定
給付企業年金の事業及び決算に関する報告書に付けて代議員会に提出し、その議決を得たうえ
で地方厚生局長等に提出しなければならない。
\end{truefalse}

\begin{answerbox}
  \textbf{誤}

  常時二十億円を下回る、又は下回ると見込まれる基金
\end{answerbox}

\newpage

\subsection{事業主等又は連合会に対する監督等}


出題例なし

\newpage

\subsection{罰則}


出題例なし

\newpage

\subsection{附則}

出題例なし

\newpage